\section{Homology review: adjunctions, split diagrams, and inverse limits}\label{note-09-homology-review-adjunctions-split-diagrams-and-inverse-limits}

The original source is the Stacks Project authors' \texttt{homology.tex} at \texttt{a04446e5\allowbreak{}7ec1fbc2\allowbreak{}52a871af\allowbreak{}cec7752f\allowbreak{}b2807b14}, SHA-256 \texttt{483CAF01\allowbreak{}B4BD164D\allowbreak{}BC9103D7\allowbreak{}6196E294\allowbreak{}E00EFE17\allowbreak{}E706178E\allowbreak{}CF2DC2C1\allowbreak{}DEF4ABE0}. This note covers original lines 7065--7945. The earlier source-bound corrections and the proposed copyedits have separate exact operations. Omitted source proofs are completed here as editorial material. No novelty is claimed, and the original statements remain identifiable.

\subsection{1. Report decisions}\label{note-09-report-decisions}

Groups 073 and 074 concern P02-E1413 and P02-E1414. The indexed objects can collectively be called a set; inserting braces around the family would clarify the notation, but the singular verb does not establish a mathematical defect or an obligatory grammatical correction. Both are recorded as optional and not integrated.

Group 075 joins P02-E1415, P02-E1416, and French HOMOLOGY-056 to MC-STK-ERR-0808. Its three existing article corrections are valid. Group 076 joins P02-E1417 and French HOMOLOGY-057 to MC-STK-ERR-0809: the sentence continues across the line break, so ``in'' takes lower case. Group 077 accepts P02-E1418, adding ``by'' at 7409. Group 078 accepts P02-E1419, adding ``by'' at 7518.

Group 079 verifies all four parenthesized colimits of MC-STK-ERR-0794, reported by French HOMOLOGY-042. Their object is the pointwise sum functor, with the exact colimit maps proved in §4. Group 080 joins P02-E1420 and French HOMOLOGY-058 to the existing article insertion MC-STK-ERR-0810. Group 081 leaves P02-E1421 unintegrated: ``the characterization of'' can identify the characterization supplied by the cited lemma. Replacing ``of'' with ``in'' is optional wording, not a mathematical correction.

Group 082 joins P02-E0017 and French HOMOLOGY-043 to MC-STK-ERR-0795: the pointwise third object is the declared \(M_i\), not an undefined \(N_i\). Group 083 verifies MC-STK-ERR-0798 from French HOMOLOGY-046: the displayed objects form a sequence of limits.

Group 084 accepts P02-E1422, adding ``by'' at 7792. Group 085 accepts P02-E1423 once, with three operations at 7815, 7816, and 7818. Group 086 joins P02-E1424 and French HOMOLOGY-059 to MC-STK-ERR-0811, which supplies the verb in the sentence introducing the short exact sequence. Group 087 verifies the degree quantifier of MC-STK-ERR-0796 from French HOMOLOGY-044. Group 088 joins P02-E0018 and French HOMOLOGY-045 to MC-STK-ERR-0797: the induction hypothesis concerns the complex \(\lim_{\gamma<\beta}K_\gamma^\bullet\), not a single degree object.

Group 089 leaves P02-E1425 unintegrated. In the ungraded three-term complex \(L_i\to M_i\to N_i\), both terms identify the same kernel-modulo-image object at the middle position. A consistent choice of ``cohomology'' is editorial style; the switch does not alter a sign, index, object, or claim. Section 8 proves the actual comparison map.

There are seventeen groups here, accounting for twenty-four reports: sixteen already corrected, four pending copyedit reports with six operations, and four optional reports. The nine existing correction units contain fourteen operations. The full Homology intake therefore has 136 reports in 89 groups, with every one of its sixty earlier units and eighty-one earlier operations reconciled.

\subsection{2. Injectives, projectives, and adjunctions}\label{note-09-injectives-projectives-and-adjunctions}

For an abelian category \(\mathcal A\), left exactness of \(\operatorname{Hom}(-,I)\) follows from the universal property of a cokernel. Surjectivity of \(\operatorname{Hom}(B,I)\to\operatorname{Hom}(A,I)\) for every monomorphism \(A\to B\) is exactly injectivity of \(I\). Hence the first two conditions in 7094--7111 agree. Injectivity splits every extension \(0\to I\to E\to C\to0\) by extending \(1_I\) to a retraction \(E\to I\). Conversely, given a monomorphism \(a:A\to B\) and \(f:A\to I\), form

\[
E=\operatorname{Coker}\bigl(A\xrightarrow{(a,-f)}B\oplus I\bigr).
\]

The induced \(I\to E\) is monic: if a morphism into \(I\) has zero image in the cokernel, the corresponding morphism into \(B\oplus I\) factors through \((a,-f)\); its \(B\)-component is zero, so monicity of \(a\) makes that factor zero. The cokernel is \(\operatorname{Coker}a\). A splitting of this extension gives a retraction \(r:E\to I\). The composite \(B\to E\xrightarrow r I\) extends \(f\), since the defining relation identifies \((a(x),0)\) with \((0,f(x))\). Thus the splitting condition implies injectivity. In the source's extension-group interpretation of \(\operatorname{Ext}\), its zero element is the split extension; this gives the fourth equivalence.

For projectives, \(\operatorname{Hom}(P,-)\) is left exact by the kernel universal property; lifting through every epimorphism is exactness. Projectivity splits every epimorphism onto \(P\). Conversely pull back an epimorphism \(A\to B\) along \(P\to B\). The pullback \(A\times_B P\to P\) is epic, and a section followed by the projection to \(A\) is the requested lift. This proves the dual extension-group criterion, retaining the original positions of all objects.

For \(J=\prod_\omega I_\omega\), extend the composite to each \(I_\omega\) across a monomorphism and combine the extensions by the product universal property. Their combined map extends the original map to \(J\); thus \(J\) is injective whenever the product exists. For \(P=\coprod_\omega P_\omega\), lift the restriction to every summand through a given epimorphism and combine the lifts by the coproduct universal property. This proves projectivity whenever that coproduct exists. The choices here are the usual set-indexed choices of lifts; no exactness of arbitrary products in \(\mathcal A\) is used.

Now let \(v:\mathcal B\to\mathcal A\) be left adjoint to \(u:\mathcal A\to\mathcal B\), with the original additive hypotheses. A left adjoint preserves cokernels, so preservation of monomorphisms is equivalent to exactness of \(v\). A lifting problem into \(uI\) is, under adjunction, a lifting problem into \(I\) along the image under \(v\) of the same monomorphism. This proves preservation of injectives when \(v\) is exact. For the converse, embed \(A=\ker(vf)\) in an injective \(I\) and extend this embedding to \(vB\). Its adjoint \(B\to uI\) extends across the monomorphism \(f:B\to B'\). Taking the adjoint back gives \(k:vB'\to I\) whose composite with \(vf\) is the chosen extension. Restricted to \(A\), it is both the original embedding and zero. Hence \(A=0\).

If \(v\) is exact and detects zero objects, an embedding \(a:vB\to I\) has monic adjoint \(j:B\to uI\). To see this, let \(e:K\to B\) be its kernel. Naturality of adjunction identifies \(je=0\) with \(a\,v(e)=0\). Both \(a\) and \(v(e)\) are monic, so \(vK=0\), and zero detection gives \(K=0\). This supplies the omitted kernel argument in the enough-injectives lemma. For an exact functor, detecting zero objects is equivalent to faithfulness: if \(v(f)=0\), exactness identifies \(v(\operatorname{im}f)\) with zero, so \(\operatorname{im}f=0\), whence \(f=0\). Conversely faithfulness applied to \(1_B\) proves zero detection. Functorial embeddings are obtained by the explicit map \(B\xrightarrow{\eta_B}uvB\to uJ(vB)\); naturality of the unit and of the chosen embedding gives its naturality, and the same kernel argument gives monicity. The already assumed functorial embeddings imply enough injectives; that redundant assumption causes no defect.

\subsection{3. The partially defined left adjoint}\label{note-09-the-partially-defined-left-adjoint}

In 7387--7454 the functor \(u\) is initially arbitrary. It is not necessary to insert an additive-functor assumption. For each generator \(P\in\mathcal P\), the representing bijection evaluated at the zero object shows that \(\operatorname{Hom}(P,u0)\) is a singleton. Choose an epimorphism from an element of \(\mathcal P\) onto \(u0\). That epimorphism must be zero, so \(u0=0\). Thus \(u\) preserves zero morphisms, and naturality shows that all the representing bijections preserve their distinguished zero elements.

Fix \(P_2\xrightarrow d P_1\xrightarrow p B\to0\). The map \(v(d)\) is uniquely determined by \(u(v(d))i_{P_2}=i_{P_1}d\). This rule also proves the identity and composition laws on the full subcategory \(\mathcal P\). Put \(v(B)=\operatorname{Coker}(v(d))\), with quotient \(q\). A morphism \(h:v(P_1)\to A\) factors uniquely through \(q\) if and only if \(h v(d)=0\). Under the representing bijection this condition is precisely \((u(h)i_{P_1})d=0\), which is equivalent to a unique factorization through \(p\). Therefore

\[
\operatorname{Hom}_{\mathcal A}(v(B),A)
\simeq \operatorname{Hom}_{\mathcal B}(B,uA)
\]

naturally in \(A\). In particular the comparison written using kernels in the source is valid as a bijection of zero fibers, without silently presupposing additivity of \(u\). For \(t:B\to B'\), precomposition gives a natural map from the functor represented by \(v(B')\) to the one represented by \(v(B)\). Yoneda supplies the unique \(v(t):v(B)\to v(B')\). Uniqueness proves composition and independence of the selected presentations up to the unique compatible isomorphism. These are the exact maps giving the left adjoint.

\subsection{4. Colimits of direct sums and the precise splitting criterion}\label{note-09-colimits-of-direct-sums-and-the-precise-splitting-criterion}

Write \(H_i=F_i\oplus G_i\), with inclusions \(a_i,b_i\) and projections \(p_i,q_i\). If \(W=\operatorname{colim}H\) exists, write its structure maps as \(c_i:H_i\to W\). The natural idempotents \(a_i p_i\) induce \(e:W\to W\), characterized by \(e c_i=c_i a_i p_i\). Equality after all \(c_i\) determines a map out of \(W\), so \(e^2=e\).

Suppose \(e=s r\), where \(r:W\to X\), \(s:X\to W\), and \(rs=1_X\). Define \(x_i=r c_i a_i:F_i\to X\). For any cocone \(f_i:F_i\to Z\), the maps \(f_i p_i:H_i\to Z\) induce a unique \(h:W\to Z\). They satisfy \(he=h\). The map \(hs:X\to Z\) then obeys \(hs x_i=f_i\). Conversely any map with that property has composite with \(r\) equal to \(h\), since \(r c_i=r e c_i=r c_i a_i p_i=x_i p_i\). Its composite with \(rs=1_X\) is therefore \(hs\), proving uniqueness. Thus \(X\) is the colimit of \(F\).

Conversely, if \(X=\operatorname{colim}F\) exists, the cocone \(c_i a_i\) induces \(s:X\to W\), while the maps \(x_i p_i:H_i\to X\) induce \(r:W\to X\). Checking on all \(x_i\) and all \(c_i\) gives \(rs=1_X\) and \(sr=e\). We have proved the exact editorial strengthening HOMOLOGY-UNDER-010:

\begin{itemize}
\tightlist
\item
  given \(W=\operatorname{colim}(F\oplus G)\), the colimit of \(F\) exists if and only if this particular \(e\) splits;
\item
  the colimit of \(G\) exists if and only if \(1-e\) splits;
\item
  both colimits exist if and only if both idempotents split, in which case their inclusions and projections give \(W\simeq\operatorname{colim}F\oplus\operatorname{colim}G\).
\end{itemize}

The last isomorphism follows because the two composite endomorphisms on \(W\) are \(e\) and \(1-e\), their sum is \(1_W\), and their cross composites are zero. Conversely, if the two colimits exist, their biproduct has the colimit universal property: a cocone on \(F\oplus G\) is exactly a pair of cocones on \(F\) and \(G\). The Karoubian assumption in the original lemma ensures both required splittings. Splitting one idempotent has not been asserted to split its complement in an arbitrary additive category. This finding specializes the earlier idempotent and retract calculations in HOMOLOGY-UNDER-001 and HOMOLOGY-UNDER-002 to a specific receiving colimit, without changing those earlier results.

For essential constancy, use the precise definition in \texttt{categories.tex} 2930--2961 at the pinned current commit. If \(H\) is essentially constant with value \(W=X\oplus Y\), let \(t:W\to H_i\) be its witness. Then \(p_i t s_X:X\to F_i\) is a witness for \(F\): its composite to \(X\) is the identity, and restricting the equation \(H_j\to H_k=(H_i\to H_k)t c_j\) to \(F_j\), then projecting to \(F_k\), gives the required factorization through \(X\). The same construction with \(q_i\) and \(s_Y\) works for \(G\). Conversely, take witnesses for \(F\) and \(G\), move them to a common index, and combine them by direct sum. For any index \(j\), the two factorization equalities hold at two later indices. Filteredness gives a common later index; its equalization axiom lets us equalize the finitely many parallel arrows from the witness index and from \(j\). After those postcompositions, both equalities hold along the same two index arrows, so their direct sum is the defining equality for \(H\). Thus the source equivalence holds. More precisely, its forward direction only requires the two specified idempotents on its value to split. The backward direction needs no Karoubian hypothesis.

For the preceding decomposition criterion in 7466--7530, a witness \(X\to M_i\) supplies a split section at every index reached from \(i\). Split the associated idempotent there; its complementary summand is \(\ker(M_j\to X)\). The defining eventual factorization kills that summand. Conversely, after choosing a split section at one index, kill the complementary summand of \(M_j\) at a later common index \(k\). The difference between the two candidate maps \(M_j\to M_k\) has zero projection to \(X\), so it factors through the complementary summand of \(M_k\). Kill this at a still later index. The difference is then zero, which is exactly the defining factorization. Applying the same argument in the opposite category gives the inverse-diagram assertion with all arrows reversed.

\subsection{5. Exact sequences and the Mittag--Leffler condition}\label{note-09-exact-sequences-and-the-mittagleffler-condition}

Kernels, cokernels, addition, zero objects, and biproducts of inverse systems are computed at each index. For example, the transition on \(\ker(L_i\to M_i)\) is the unique map whose composite with the kernel inclusion is the restriction of the transition on \(L_i\). The kernel universal property gives its composition law and its universal factorization for maps of systems. The cokernel construction is dual. The canonical coimage-to-image map is an isomorphism at each index, and its inverses commute with transitions by uniqueness. This proves the additive and abelian assertions and their pointwise exactness criterion, including the corrected object \(M_i\) at 7624.

For a short exact sequence \(0\to(A_i)\to(B_i)\to(C_i)\to0\) of abelian groups, a compatible tuple in \(B_i\) mapping to zero in each \(C_i\) comes from a unique compatible tuple in \(A_i\). This proves left exactness of inverse limits. If \((B_i)\) is ML, \(\operatorname{im}(C_k\to C_i)\) is the image in \(C_i\) of \(\operatorname{im}(B_k\to B_i)\); hence these images stabilize too.

Suppose now \((A_i)\) is ML and fix \(c=(c_i)\in\lim C_i\). The nonempty fibers \(T_i\subset B_i\) over \(c_i\) form an inverse system of sets. Each image of \(T_k\to T_i\) is a nonempty coset of \(\operatorname{im}(A_k\to A_i)\). These cosets are nested, and once those subgroups stabilize, the nonempty nested cosets are equal. Let \(U_i\subset T_i\) be that stable image. Its transition to \(U_{i-1}\) is surjective: given a member of \(U_{i-1}\), lift it from \(T_k\) with \(k\) beyond stabilization at both indices; its image at \(i\) belongs to \(U_i\). Choose successively compatible members of the nonempty \(U_i\). They give a lift of \(c\) in \(\lim B_i\). This proves the exactness assertion used below.

The lack of general right exactness can be seen without dropping any transition data. Set \(A_i=B_i=\mathbf Z\), embed \(A_i\) into \(B_i\) by multiplication by \(2^i\), use multiplication by 2 on \(A_{i+1}\to A_i\), and the identity on \(B_{i+1}\to B_i\). Then \(C_i=\mathbf Z/2^i\mathbf Z\) has reduction transitions. Let \(c_i\) be the residue of \(\sum_{0\leq 2k<i}2^{2k}\). These residues are compatible and \(3c_i=-1\) modulo \(2^i\). If they came from an integer \(b\), then \(3b+1\) would be divisible by every \(2^i\), forcing \(3b=-1\), which is impossible in \(\mathbf Z\).

For the four-term exact sequence at 7700, put \(I_i=\operatorname{im}(A_i\to B_i)\) and \(Z_i=\ker(C_i\to D_i)\). The quotient argument makes \((I_i)\) ML, and \(0\to I_i\to B_i\to Z_i\to0\) gives surjectivity \(\lim B_i\to\lim Z_i\). Left exactness identifies \(\lim Z_i\) with the required kernel, proving the assertion.

An essentially constant inverse system over \(\mathbf N\) has a retraction \(A_i\to A\) at some index. For \(j\geq i\), compose \(A_j\to A_i\to A\). Its section is the limit projection \(A\to A_j\), so \(A_j=A\oplus Z_j\), compatibly with transitions. The defining eventual factorization kills each fixed \(Z_j\) after a sufficiently late transition. Conversely these decompositions give that factorization by projecting to \(A\). In particular the stable image at each fixed index is the image of \(A\), proving ML.

For 7760--7803, first take an essentially zero quotient system \((Z_i)\). Choose \(n_1<n_2<\cdots\), cofinal in \(\mathbf N\), with \(Z_{n_{j+1}}\to Z_{n_j}\) zero. For fixed \(i\) and \(j>i\), the source's exact inclusions are

\[
\operatorname{im}(A_{n_j}\to A_{n_i})
\subset\operatorname{im}(B_{n_j}\to B_{n_i})
\subset\operatorname{im}(A_{n_{j-1}}\to A_{n_i}).
\]

Stabilization of the outer images forces stabilization of the middle ones. Conversely the middle images for \(j+1\) and \(j\) enclose the first image for \(j\), proving the reverse implication. A cofinal subsequence detects ML: for any original index, map from a later subsequence index and squeeze every subsequent image between two subsequence images. Thus the conclusion holds for the full systems. For a quotient \(C\oplus Z_i\), replace \(B_i\) by the inverse image \(B'_i\) of \(C\); the first case handles \(B/B'=Z\). With constant quotient \(C\), there is an exact sequence

\[
0\to\operatorname{im}(A_j\to A_i)
\to\operatorname{im}(B_j\to B_i)\to C\to0.
\]

To verify its kernel, a lift from \(B_j\) mapping to zero in \(C\) already belongs to \(A_j\), because the transition on \(C\) is the identity. Nested image subgroups with the same kernel and with surjective map to \(C\) are equal: lift an element's image in \(C\) from the smaller subgroup and subtract; the difference belongs to the common kernel. This proves both directions of stabilization.

\subsection{6. Cohomology and countable limits}\label{note-09-cohomology-and-countable-limits}

Retain the four original terms \(A_i^{-2}\to A_i^{-1}\to A_i^0\to A_i^1\), and write \(Z_i^j=\ker d_i^j\), \(I_i^j=\operatorname{im}d_i^{j-1}\). Under the stated hypotheses, \((I_i^{-1})\) and \((I_i^0)\) are ML as quotients of \((A_i^{-2})\) and \((A_i^{-1})\). The exact sequence \(0\to I_i^{-1}\to Z_i^{-1}\to H_i^{-1}\to0\), with essentially constant last system, then makes \((Z_i^{-1})\) ML by §5. Applying §5 to \(0\to Z_i^{-1}\to A_i^{-1}\to I_i^0\to0\) gives \(\lim A_i^{-1}\twoheadrightarrow\lim I_i^0\). Applying it to \(0\to I_i^0\to Z_i^0\to H_i^0\to0\) gives

\[
0\to\lim I_i^0\to\lim Z_i^0\to\lim H_i^0\to0.
\]

The kernel universal property identifies \(\lim Z_i^0\) with \(\ker(\lim A_i^0\to\lim A_i^1)\), and the first surjection identifies \(\lim I_i^0\) with the image of \(\lim A_i^{-1}\to\lim A_i^0\). Taking the quotient proves the source's isomorphism. Explicitly it sends the class of a compatible cycle tuple \((x_i)\) to \(([x_i])\). The two exact sequences prove both injectivity and surjectivity of this particular map.

\subsection{7. Ordinal limits and a bounded-degree consequence}\label{note-09-ordinal-limits-and-a-bounded-degree-consequence}

Let \(r_\beta:K_\beta^\bullet\to L_\beta^\bullet\), where \(L_\beta^\bullet=\lim_{\gamma<\beta}K_\gamma^\bullet\), be the matching map. The limit over the empty ordinal is the zero complex. For a fixed integer \(N\), suppose

\[
H^q(K_\beta^\bullet)=0\quad(q\leq N),\qquad
r_\beta^j\text{ is surjective}\quad(j\leq N-2)
\]

for every \(\beta<\alpha\). Then \(H^q(\lim_{\beta<\alpha}K_\beta^\bullet)=0\) for all \(q\leq N\). This is HOMOLOGY-UNDER-011. Its hypotheses are actual data about the original complexes and matching maps; no vanishing of a limit is assumed.

Prove the assertion by transfinite induction on \(\alpha\), retaining all degrees \(q\leq N\) simultaneously. It holds for \(\alpha=0\). For a cycle \(x=(x_\beta)\) of degree \(q\), construct compatible primitives \(y_\beta\in K_\beta^{q-1}\) recursively. At stage \(\beta\), choose \(y'_\beta\) with \(dy'_\beta=x_\beta\), using \(H^q(K_\beta)=0\). The previously chosen tuple \(y_{<\beta}\) belongs to \(L_\beta^{q-1}\). Its difference from \(r_\beta(y'_\beta)\) is a cycle, since both differentials equal \(r_\beta(x_\beta)\). The induction hypothesis at \(\beta<\alpha\) gives \(H^{q-1}(L_\beta)=0\). Consequently

\[
y_{<\beta}-r_\beta(y'_\beta)=dz
\quad\text{for some }z\in L_\beta^{q-2}.
\]

Surjectivity of \(r_\beta^{q-2}\), since \(q-2\leq N-2\), gives \(\widetilde z\in K_\beta^{q-2}\) with \(r_\beta\widetilde z=z\). Set \(y_\beta=y'_\beta+d\widetilde z\). Its differential remains \(x_\beta\), and its entire tuple of preceding images is \(y_{<\beta}\). At \(\beta=0\), that tuple and \(z\) are zero. The recursion therefore works at both successor and limit ordinals and gives a compatible \(y\) with \(dy=x\). This proves the claim.

The original acyclicity assertion follows by applying this result for every integer \(N\), because the source assumes acyclicity in every degree and matching surjectivity in every degree. This verifies both the inserted quantifier and the corrected complex-valued induction hypothesis. The stronger bounded-degree statement does not require matching surjectivity in degrees \(N-1,N,N+1,\ldots\), nor any cohomology vanishing above \(N\). It is retained as an editorial consequence, not substituted for the source lemma.

\subsection{8. Exactness of products}\label{note-09-exactness-of-products}

For the source's complexes \(L_i\xrightarrow{a_i}M_i\xrightarrow{b_i}N_i\), the product kernel is exactly \(\prod_i\ker b_i\): the equation \((b_i(m_i))=0\) is equivalent to all its coordinate equations. The image of \(\prod_i a_i\) is exactly \(\prod_i\operatorname{im}a_i\): one inclusion follows coordinatewise, and the other follows by choosing a preimage of each coordinate. The map

\[
\frac{\prod_i\ker b_i}{\prod_i\operatorname{im}a_i}
\longrightarrow \prod_i(\ker b_i/\operatorname{im}a_i),
\qquad [(m_i)]\longmapsto([m_i])
\]

is injective because its zero fiber consists of exactly the indicated product image, and it is surjective by choosing a representative of each coordinate class. For an empty index set both groups are zero, so the same statement holds. This establishes the original comparison without making a chain-versus-cochain convention out of the ungraded three-term notation.

\subsection{9. Propagation and bounded reading scope}\label{note-09-propagation-and-bounded-reading-scope}

The direct split-diagram use in \texttt{derived.tex} 5168--5232 invokes the forward essential-constancy assertion to separate the diagrams of \(X\) and \(Y\). Its stated Karoubian hypothesis supplies both splittings, so that use remains valid. The exact local requirement exposed in §4 is the splitting of the two projection idempotents on the colimit value. It does not by itself certify the separate cofinality argument later in that receiving proof.

The ordinal-limit uses at \texttt{injectives.tex} 1700 and \texttt{sdga.tex} 3962 invoke the original all-degree assertion on complexes of homomorphisms. Their displayed constructions use surjective restrictions at successor stages and equality with the earlier limit at limit stages. Thus their matching-map requirement has the correct form, and the original Homology implication is preserved. These bounded reads do not certify the entire transfinite constructions or unrelated SDGA formulas.

The countable cohomology comparison is cited in \texttt{more-algebra.tex} 6504 and 6526 and in \texttt{cohomology.tex} 10726. At the first More Algebra use, the degree -2 and -1 terms are finite free modules tensored with the given surjective system, hence have surjective transitions; the preceding Tor system is identified there as essentially zero. At the second use, both Tor systems are essentially zero and therefore ML; the four-term sequence is exact and its degree -1 cohomology is zero. These supply the exact Homology hypotheses. The Tor comparison and Artin--Rees input are cited receiving data, not newly certified results. The Cohomology use explicitly has eventual constancy of cohomology; the earlier construction supplying its termwise surjectivity remains outside the displayed bounded interval, so that use is routed without a claim of complete verification.

The link at \texttt{more-algebra.tex} 23643 points to that chapter's stronger treatment. Its complete statement and proof have not been read in this batch; the bounded-degree consequence in §7 has no novelty claim. Product-exactness references in Cohomology and Sites Cohomology are listed in the receiving-route ledger; no mathematical change arose from the optional terminology report, and no unperformed review of those receiving passages is claimed.

Actual consecutive reading now covers all 7,945 original source lines. All received Homology reports have decisions. Candidate preparation, integration, builds, rendered-page inspection, and publication remain subsequent work; this proof note is not an admission or release receipt.
